\scanpage{142}
\mspage{91}

\diagram{\((**)\qquad\begin{tikzcd}[ampersand replacement=\&]
 \& V \arrow[dl,"{\widetilde f=f_{/e_V}}"'] \arrow[dd,"f","{\deleted{\unsure{V'|V}}}"{pos=.65,description}]\\
 \begin{gathered}\underline{\pi}_0(V/V')\\=V'_{/f_!(e_V)}\end{gathered} \arrow[dr,"\mathrm{can}"'] \& \\
 \& V'
\end{tikzcd}\)}.

Explicitons le $V'$-morphisme $\widetilde f$ sur les
fibres,
\[
 \widetilde f_{v'}:V_{v'}\longrightarrow
 \underline{\pi}_0(V/\unsure{v'|V'})_{v'}\overset{\text{déf}}{=}\pi_0({}_{v'}\backslash V),\tag*{(***)}
\]
où \inserted{dans} le \deleted{\unsure{deuxième}}\inserted{dernier} membre, on a identifié
l'ensemble $\pi_0({}_{v'}\backslash V)$ à la catégorie discrète
qu'il définit. Pour expliciter cette flèche
$(***)$, notons qu'on a une inclusion canonique
\[
 V_{v'}\longrightarrow{}_{v'}\backslash V
\]
et pour toute catégorie $W$, un morphisme
canonique (d'adjonction) $W\unsure{-|\to}\pi_0(W)$,
d'où un composé
\[
 V_{v'}-{}_{v'}\backslash V-\pi_0({}_{v'}\backslash V),
\]
qui n'est autre que $\widetilde f_{v'}$ dans $(***)$.

Il doit être facile de montrer que
la factorisation $(**)$ de $f$ est caractérisée,
à isom. \unsure{unique|uniques} près, \deleted{\unsure{par les}} en tant que

% archive sheet number: 142
